\section{Pruebas de hipótesis}
Una prueba de hipótesis es un procedimiento estadístico estándar para tomar una decisión
objetiva entre dos afirmaciones opuestas sobre una población. Sus componentes básicos son:
\begin{itemize}
  \item \textbf{La hipótesis nula ($H_0$):} es la afirmación que se asume como verdadera y se    somete a prueba. Siempre incluye el signo igual.
  \item \textbf{La hipótesis alternativa ($H_1$):} es la sospecha o el resultado deseado por el    investigador. Determina hacia qué lado de la campana se ubica la zona de rechazo.
\end{itemize}
Para tomar la decisión, el investigador define un nivel de significación ($\alpha$), usualmente
del 5\%, que delimita la zona de rechazo mediante un valor crítico ($Z_c$). Si el estadístico
de prueba ($Z_0$) calculado con los datos muestrales cae dentro de esa región crítica, la
hipótesis nula se rechaza.

\AnalisisDescriptivo
Para poner en práctica la teoría tomamos el siguiente problema: la cátedra de Física afirma que
el puntaje promedio histórico de los alumnos es de 82 puntos. Evaluamos una muestra aleatoria
de estudiantes para contrastar esta afirmación frente a la sospecha de que el promedio real es
mayor. Los datos del problema son:

\begin{TablaDatosFlexible}{L{0.30\linewidth} C{0.18\linewidth} L{0.44\linewidth}}{\textbf{Parámetro} & \textbf{Valor} & \textbf{Significado práctico}}

\TablaFilaFlexible{Media histórica ($\mu_0$) & 82 & El puntaje promedio que la cátedra de Física afirma como histórico.}
\TablaFilaFlexible{Cantidad de datos ($n$) & 50 & Estudiantes de la muestra aleatoria seleccionada para el contraste.}
\TablaFilaFlexible{Media muestral ($\bar{x}$) & 83,46 & El promedio de física obtenido por esa muestra de 50 alumnos.}
\TablaFilaFlexible{Desviación estándar ($\sigma$) & 3,5 & Dispersión de los puntajes de física, conocida para la población.}
\TablaFilaFlexible{Nivel de significación ($\alpha$) & 0,05 & Probabilidad de equivocarnos al rechazar $H_0$ (5\%).}
\TablaFilaFlexible{Valor crítico ($Z_c$) & 1,645 & Umbral de la zona de rechazo para una prueba unilateral derecha al 5\%.}

\end{TablaDatosFlexible}

\subsection{Planteo estadístico del problema}
\begin{itemize}
  \item \textbf{Hipótesis nula:} $H_0: \mu \leq 82$ (afirmación histórica de la cátedra).
  \item \textbf{Hipótesis alternativa:} $H_1: \mu > 82$ (sospecha del investigador — prueba
    unilateral derecha).
  \item \textbf{Datos muestrales:} media calculada $\bar{x} = 83{,}46$ puntos y desviación
    estándar $\sigma = 3{,}5$ puntos, con $n = 50$ estudiantes.
  \item \textbf{Criterio de decisión:} con $\alpha = 0{,}05$ (5\%), el valor crítico es
    $Z_c = 1{,}645$.
\end{itemize}

\subsection{Cálculo del estadístico de prueba ($Z_0$)}
\begin{equation*}
  Z_0 = \frac{\bar{x} - \mu_0}{\dfrac{\sigma}{\sqrt{n}}}
      = \frac{83{,}46 - 82}{\dfrac{3{,}5}{\sqrt{50}}}
      = \frac{1{,}46}{0{,}495} = \mathbf{2{,}95}
\end{equation*}

% --------------------------------------------------------------------------------------------------

\subsection{Evaluación y conclusión}
\PalabraClave{Regla de decisión:} Si $Z_0 > Z_c$, se rechaza $H_0$.

\PalabraClave{Resultado:} Dado que el estadístico obtenido ($2{,}95$) es mayor que el valor
crítico ($1{,}645$), el dato cae profundamente dentro de la zona de rechazo.

\PalabraClave{Conclusión del estudio:} Con un nivel de significación del 5\%, existe evidencia
estadística suficiente para rechazar la hipótesis nula. Concluimos que el puntaje promedio en
física de los estudiantes es significativamente superior a los 82 puntos históricos, validando
por completo nuestra sospecha inicial.

\TituloGrafico{Zona de rechazo y estadístico de prueba}
\begin{figure}[H]
    \centering
    \begin{tikzpicture}
        % Función de densidad de la normal estándar.
        \pgfmathdeclarefunction{normalpdf}{1}{%
          \pgfmathparse{1/(sqrt(2*pi))*exp(-(#1)^2/2)}%
        }
        \begin{axis}[
            width=0.66\textwidth,
            height=6.2cm,
            xlabel={Estatístico ($Z$)},
            ylabel={},
            xmin=-3.6, xmax=3.8,
            ymin=0, ymax=0.46,
            xtick={-3,-2,-1,0,1,2,3},
            ytick=\empty,
            axis lines=left,
            samples=200,
            legend style={at={(1.05,1.0)}, anchor=north west, font=\small,
                fill=white, draw=black!50, fill opacity=1, text opacity=1,
                inner sep=4pt, cells={anchor=west}},
            legend cell align=left,
            clip=false,
        ]
            % Zona de rechazo: Z > 1,645 (5% de la cola derecha).
            \addplot[fill=red!70!black, draw=none, fill opacity=0.25, domain=1.645:3.8]
                {normalpdf(x)} \closedcycle;
            \addlegendentry{Zona de rechazo ($\alpha = 0{,}05$)}
            % Campana normal estándar.
            \addplot[blue, thick, domain=-3.6:3.8] {normalpdf(x)};
            \addlegendentry{Normal estándar ($H_0$)}
            % Entradas de leyenda para elementos dibujados con \draw.
            \addlegendimage{red!70!black, dashed, thick}
            \addlegendentry{Valor crítico ($Z_c$)}
            \addlegendimage{black, thick, -{Latex[length=2.5mm]}}
            \addlegendentry{Estadístico ($Z_0$)}
            % Valor crítico.
            \draw[red!70!black, dashed, thick] (axis cs:1.645,0) -- (axis cs:1.645,0.415)
                node[above, red!70!black] {$Z_c = 1{,}645$};
            % Estadístico de prueba con flecha desde el valor crítico.
            \draw[black, thick, -{Latex[length=2.5mm]}] (axis cs:1.645,0.10)
                -- (axis cs:2.95,0.10);
            \node[above, black] at (axis cs:2.95,0.10) {$Z_0 = 2{,}95$};
            % Etiqueta de la zona de rechazo.
            \node[red!70!black] at (axis cs:2.7,0.30) {rechazo $H_0$};
        \end{axis}
    \end{tikzpicture}
    \caption{Distribución normal estándar con la zona de rechazo de una prueba unilateral
    derecha al 5\% (área roja a la derecha de $Z_c = 1{,}645$). El estadístico $Z_0 = 2{,}95$
    cae dentro de esa zona, por lo que se rechaza $H_0$.}
    \label{fig:hipotesis_rechazo}
\end{figure}

\Interpretacion{La prueba muestra que la diferencia entre el promedio muestral (83,46) y el histórico (82) es demasiado grande para atribuirse al azar: con $Z_0 = 2{,}95$ muy por encima de $Z_c = 1{,}645$, la probabilidad de obtener esa diferencia si $H_0$ fuera cierta es inferior al 5\%. Por eso rechazamos la hipótesis nula y aceptamos la sospecha: el puntaje promedio en física es significativamente superior a los 82 puntos históricos. Vale recordar que $\alpha = 0{,}05$ representa el riesgo que asumimos de equivocarnos al tomar esta decisión.}